\clearpage\thispagestyle{empty}
\pdfbookmark[0]{Sources for the review of Connoissance des tems for 1812}{connoissance1812-sources}
\begin{center}\textbf{Sources for the review of Connoissance des tems for 1812}\end{center}
\noindent This AI-translated review follows the editable transcription of Carl Friedrich Gauss, \textit{Werke}, vol.~VI (Göttingen, 1874), pp.~537--541. It first appeared under Paris in \textit{Göttingische gelehrte Anzeigen}, no.~199, 14 December 1811, pp.~1978--1987. The reviewed \href{https://cdt.imcce.fr/items/show/161}{almanac for 1812} is catalogued by the Institut de mécanique céleste et de calcul des éphémérides.
\par\smallskip\noindent The date after \textit{Juillet} appears as 181c in the collected printing and as 1310 in the digitized first printing. The translation retains 181c. The collected scan reads \textit{von Gauss} in the first Delambre heading and appears to read \textit{carporum} in the title of the \textit{Theoria motus}. The editable transcription reads \textit{vom Gauss} and \textit{corporum}.
\par\smallskip\noindent The review cites the September 1808 issue of the \textit{Monatliche Correspondenz}. The related article in \textit{Werke}, vol.~VI, p.~129, gives October 1808 in its heading. The last term of the latitude-error expression on p.~139 of that article has a minus sign. The review gives a plus sign.
\par\smallskip\noindent The first French quotation says that Burckhardt compared and calculated afresh \emph{all} the observations, old and modern alike. Delambre's first remark says the last formula proves that in practice the method would have only fairly limited accuracy. His second says that on a night when three stars could be observed at the same altitude, one would very probably see a known star cross the meridian and obtain the latitude with less effort and greater precision, and the time shown by the pendulum clock from a single altitude.
\par\smallskip\noindent \href{https://gdz.sub.uni-goettingen.de/id/PPN236007467}{Collected text}. First printing, \href{https://images.p.sub.uni-goettingen.de/iiif/image/gdz:PPN319721507_1811:00002012/full/1650,/0/default.jpg}{p.~1978}, \href{https://images.p.sub.uni-goettingen.de/iiif/image/gdz:PPN319721507_1811:00002013/full/1650,/0/default.jpg}{p.~1979}, \href{https://images.p.sub.uni-goettingen.de/iiif/image/gdz:PPN319721507_1811:00002014/full/1650,/0/default.jpg}{p.~1980}, \href{https://images.p.sub.uni-goettingen.de/iiif/image/gdz:PPN319721507_1811:00002015/full/1650,/0/default.jpg}{p.~1981}, \href{https://images.p.sub.uni-goettingen.de/iiif/image/gdz:PPN319721507_1811:00002016/full/1650,/0/default.jpg}{p.~1982}, \href{https://images.p.sub.uni-goettingen.de/iiif/image/gdz:PPN319721507_1811:00002017/full/1650,/0/default.jpg}{p.~1983}, \href{https://images.p.sub.uni-goettingen.de/iiif/image/gdz:PPN319721507_1811:00002018/full/1650,/0/default.jpg}{p.~1984}, \href{https://images.p.sub.uni-goettingen.de/iiif/image/gdz:PPN319721507_1811:00002019/full/1650,/0/default.jpg}{p.~1985}, \href{https://images.p.sub.uni-goettingen.de/iiif/image/gdz:PPN319721507_1811:00002020/full/1650,/0/default.jpg}{p.~1986} and \href{https://images.p.sub.uni-goettingen.de/iiif/image/gdz:PPN319721507_1811:00002021/full/1650,/0/default.jpg}{p.~1987}.
\par\smallskip\noindent Translation and typesetting produced by AI.
\clearpage
